\documentclass[11pt, a4paper]{article}

\usepackage[utf8]{inputenc}
\usepackage[T1]{fontenc}
\usepackage[spanish]{babel}
\usepackage{geometry}
\usepackage{amsmath}
\usepackage{booktabs}
\usepackage{listings}
\usepackage{xcolor}
\usepackage{hyperref}
\usepackage{parskip}
\usepackage{titlesec}
\usepackage{enumitem}
\usepackage{fancyhdr}
\usepackage{graphicx}

\geometry{margin=2.5cm}

% Colores para código
\definecolor{codebg}{RGB}{245,245,245}
\definecolor{codeframe}{RGB}{200,200,200}
\definecolor{keyword}{RGB}{0,0,180}
\definecolor{comment}{RGB}{100,100,100}
\definecolor{string}{RGB}{180,0,0}

\lstset{
  backgroundcolor=\color{codebg},
  frame=single,
  rulecolor=\color{codeframe},
  basicstyle=\ttfamily\small,
  keywordstyle=\color{keyword}\bfseries,
  commentstyle=\color{comment}\itshape,
  stringstyle=\color{string},
  breaklines=true,
  numbers=left,
  numberstyle=\tiny\color{gray},
  numbersep=8pt,
  showstringspaces=false,
  language=Python,
}

\lstdefinestyle{bash}{
  language=bash,
  backgroundcolor=\color{black!90},
  basicstyle=\ttfamily\small\color{white},
  keywordstyle=\color{green!70},
  frame=none,
  numbers=none,
}

\pagestyle{fancy}
\fancyhf{}
\rhead{Catan AI -- Documentación Técnica}
\lhead{Fase 3: Red Neuronal}
\rfoot{Página \thepage}

\hypersetup{
  colorlinks=true,
  linkcolor=blue!60!black,
  urlcolor=blue!60!black,
}

% ──────────────────────────────────────────────────────────────────────────────

\begin{document}

\begin{titlepage}
  \centering
  \vspace*{3cm}
  {\Huge\bfseries Catan AI\par}
  \vspace{0.5cm}
  {\Large Documentación Técnica -- Fase 3\par}
  \vspace{0.3cm}
  {\large Red Neuronal con Aprendizaje por Refuerzo\par}
  \vspace{2cm}
  \rule{0.6\textwidth}{0.4pt}
  \vspace{0.5cm}

  \begin{tabular}{ll}
    \textbf{Entorno:}    & \texttt{catanatron}       \\
    \textbf{Red:}        & PyTorch (CPU)             \\
    \textbf{Algoritmo:}  & REINFORCE                 \\
    \textbf{Estado:}     & Vector de 250 features    \\
    \textbf{Acciones:}   & Espacio fijo de 244       \\
  \end{tabular}

  \vspace{2cm}
  \rule{0.6\textwidth}{0.4pt}
  \vspace{0.5cm}

  {\large Proyecto desarrollado en Python}\par
  \vspace{0.3cm}
  {\large\textit{Stack: catanatron + PyTorch + matplotlib}}\par
  \vfill
\end{titlepage}

% ──────────────────────────────────────────────────────────────────────────────
\tableofcontents
\newpage

% ──────────────────────────────────────────────────────────────────────────────
\section{Descripción General}

Este documento describe la implementación de la \textbf{Fase 3} del proyecto Catan AI: el
diseño, entrenamiento y evaluación de una red neuronal que aprende a jugar al Catan mediante
aprendizaje por refuerzo (\textit{Reinforcement Learning}, RL).

\subsection{Benchmarks de referencia}

\begin{center}
\begin{tabular}{lcc}
  \toprule
  \textbf{Agente} & \textbf{Victorias / 200} & \textbf{Tasa de victoria} \\
  \midrule
  Aleatorio (Fase 1)   & $\sim$50 & $\sim$25\% \\
  Heurístico (Fase 2)  & 93       & 46.5\%     \\
  Red neuronal (Fase 3)& ?        & \textbf{objetivo: $>$46.5\%} \\
  \bottomrule
\end{tabular}
\end{center}

\subsection{Estructura de archivos}

\begin{center}
\begin{tabular}{ll}
  \toprule
  \textbf{Archivo} & \textbf{Propósito} \\
  \midrule
  \texttt{model.py}           & Features, codificación de acciones, arquitectura de la red \\
  \texttt{neural\_player.py}  & Jugador que usa la red para decidir \\
  \texttt{train.py}           & Bucle de entrenamiento REINFORCE \\
  \texttt{eval.py}            & Evaluación formal contra otros agentes \\
  \texttt{plot.py}            & Gráfico del progreso del entrenamiento \\
  \texttt{heuristic\_player.py} & Agente de Fase 2 (oponente de referencia) \\
  \bottomrule
\end{tabular}
\end{center}


% ──────────────────────────────────────────────────────────────────────────────
\section{Representación del Estado}

La red recibe el estado del juego como un vector numérico de dimensión $\mathbf{250}$,
construido en \texttt{model.py} por la función \texttt{extract\_features(state, color)}.

\subsection{Composición del vector}

\begin{center}
\begin{tabular}{clc}
  \toprule
  \textbf{Índices} & \textbf{Componente} & \textbf{Dimensión} \\
  \midrule
  0--99   & Estado de los 4 jugadores (25 features $\times$ 4)  & 100 \\
  100--153 & Edificios en los 54 nodos del tablero               & 54  \\
  154--225 & Caminos en las 72 aristas del tablero               & 72  \\
  226--244 & Posición del ladrón (one-hot sobre 19 fichas)       & 19  \\
  245--249 & Recursos del banco (normalizados)                   & 5   \\
  \midrule
  & \textbf{Total}                                              & \textbf{250} \\
  \bottomrule
\end{tabular}
\end{center}

\subsection{Features por jugador (25 por jugador)}

Para cada jugador, la perspectiva se \textbf{rota} de forma que el jugador RL aparece
siempre en la posición $P_0$. Esto permite que la red aprenda de forma independiente
del color asignado.

Las 25 features por jugador son:
\begin{itemize}[noitemsep]
  \item 10 stats: puntos de victoria (visibles y reales), caminos/colonias/ciudades
        disponibles, longitud de camino más largo, banderas de turno.
  \item 5 recursos en mano: WOOD, BRICK, SHEEP, WHEAT, ORE.
  \item 5 cartas de desarrollo en mano: KNIGHT, YEAR\_OF\_PLENTY, MONOPOLY,
        ROAD\_BUILDING, VICTORY\_POINT.
  \item 5 cartas de desarrollo jugadas.
\end{itemize}

\subsection{Codificación del tablero}

\begin{itemize}[noitemsep]
  \item \textbf{Nodos} (54): $+1$ colonia propia, $+2$ ciudad propia,
        $-1$ colonia rival, $-2$ ciudad rival, $0$ vacío.
  \item \textbf{Aristas} (72): $+1$ camino propio, $-1$ camino rival, $0$ sin camino.
  \item \textbf{Ladrón}: vector one-hot de 19 posiciones (fichas de tierra).
  \item \textbf{Banco}: cada recurso normalizado por 19 (cantidad máxima inicial).
\end{itemize}


% ──────────────────────────────────────────────────────────────────────────────
\section{Espacio de Acciones}

Se define un espacio de acciones \textbf{fijo} de dimensión $\mathbf{244}$, con una
máscara binaria que anula las acciones inválidas en cada estado.

\begin{center}
\begin{tabular}{clc}
  \toprule
  \textbf{Índices} & \textbf{Tipo de acción} & \textbf{Cantidad} \\
  \midrule
  0        & ROLL                              & 1  \\
  1        & END\_TURN                         & 1  \\
  2        & BUY\_DEVELOPMENT\_CARD            & 1  \\
  3        & PLAY\_KNIGHT\_CARD                & 1  \\
  4        & PLAY\_ROAD\_BUILDING              & 1  \\
  5--58    & BUILD\_SETTLEMENT (nodo 0--53)    & 54 \\
  59--112  & BUILD\_CITY (nodo 0--53)          & 54 \\
  113--184 & BUILD\_ROAD (arista 0--71)        & 72 \\
  185--203 & MOVE\_ROBBER (ficha 0--18)        & 19 \\
  204--208 & PLAY\_MONOPOLY (recurso 0--4)     & 5  \\
  209--223 & PLAY\_YEAR\_OF\_PLENTY (15 combos)& 15 \\
  224--243 & MARITIME\_TRADE (20 combos)       & 20 \\
  \midrule
  & \textbf{Total}                            & \textbf{244} \\
  \bottomrule
\end{tabular}
\end{center}

La acción \texttt{DISCARD} no se incluye en el espacio fijo por su complejidad
combinatoria; se maneja con selección aleatoria entre las opciones válidas.


% ──────────────────────────────────────────────────────────────────────────────
\section{Arquitectura de la Red Neuronal}

La red implementada en \texttt{model.py} (\texttt{CatanNet}) es un perceptrón
multicapa (\textit{Multi-Layer Perceptron}, MLP):

\[
  \text{input}\ (250) \;\to\; \underbrace{[256 \to \text{ReLU}]}_{\times 3} \;\to\; \text{output}\ (244)
\]

\begin{lstlisting}[caption={Definición de CatanNet en PyTorch}]
class CatanNet(nn.Module):
    def __init__(self, hidden_size=256, num_layers=3):
        super().__init__()
        layers = [nn.Linear(STATE_SIZE, hidden_size), nn.ReLU()]
        for _ in range(num_layers - 1):
            layers += [nn.Linear(hidden_size, hidden_size), nn.ReLU()]
        layers.append(nn.Linear(hidden_size, ACTION_SPACE_SIZE))
        self.net = nn.Sequential(*layers)

    def forward(self, x):
        return self.net(x)  # retorna logits (244,)
\end{lstlisting}

La salida son \textbf{logits} (sin softmax). La selección de acción aplica:
\begin{enumerate}
  \item Mascarar logits inválidos con $-\infty$.
  \item Aplicar softmax para obtener una distribución de probabilidad.
  \item Durante entrenamiento: muestrear (exploración estocástica).
  \item Durante evaluación: argmax (acción más probable).
\end{enumerate}


% ──────────────────────────────────────────────────────────────────────────────
\section{Algoritmo de Entrenamiento: REINFORCE}

Se utiliza el algoritmo \textbf{REINFORCE} (Policy Gradient de Monte Carlo),
implementado en \texttt{train.py}.

\subsection{Fundamento teórico}

El objetivo es maximizar la recompensa esperada $J(\theta) = \mathbb{E}_\pi[R]$.
El gradiente de la política es:

\[
  \nabla_\theta J(\theta) = \mathbb{E}_\pi \left[
    \sum_{t=0}^{T} \nabla_\theta \log \pi_\theta(a_t \mid s_t) \cdot R
  \right]
\]

donde $R \in \{+1, -1\}$ es la recompensa al final del juego (victoria o derrota).

\subsection{Función de pérdida}

\[
  \mathcal{L}(\theta) = -\frac{1}{T} \sum_{t=0}^{T} \log \pi_\theta(a_t \mid s_t) \cdot R
\]

\subsection{Hiperparámetros}

\begin{center}
\begin{tabular}{ll}
  \toprule
  \textbf{Parámetro} & \textbf{Valor} \\
  \midrule
  Optimizador        & Adam           \\
  Learning rate      & $10^{-4}$      \\
  Capas ocultas      & 3 $\times$ 256 \\
  Recompensa         & $+1$ victoria, $-1$ derrota \\
  Clipping de gradiente & $\|\nabla\| \leq 1.0$ \\
  \bottomrule
\end{tabular}
\end{center}

\subsection{Bucle de entrenamiento}

\begin{lstlisting}[caption={Episodio de entrenamiento simplificado}]
def run_episode(model, optimizer):
    rl_player = NeuralPlayer(Color.RED, model, training=True)
    others    = [HeuristicPlayer(c) for c in [BLUE, WHITE, ORANGE]]
    game = Game([rl_player] + others)
    game.play()

    reward    = +1.0 if game.winning_color() == Color.RED else -1.0
    log_probs = torch.stack(rl_player.log_probs)
    loss      = -(log_probs * reward).mean()

    optimizer.zero_grad()
    loss.backward()
    torch.nn.utils.clip_grad_norm_(model.parameters(), 1.0)
    optimizer.step()
\end{lstlisting}


% ──────────────────────────────────────────────────────────────────────────────
\section{Guía de Uso}

\subsection{Requisitos}

\begin{lstlisting}[style=bash]
cd C:\Users\franc\Documents\catan-ai
venv\Scripts\activate
\end{lstlisting}

Las dependencias necesarias son \texttt{catanatron}, \texttt{torch},
\texttt{numpy} y \texttt{matplotlib}, ya instaladas en el entorno virtual.

\subsection{Entrenamiento}

\begin{lstlisting}[style=bash]
# Entrenamiento básico (5000 episodios vs heurístico)
python train.py

# Empezar vs aleatorio (más rápido, bueno para las primeras épocas)
python train.py --episodes 2000 --opponents random --save model_pre.pt

# Continuar entrenamiento desde un checkpoint
python train.py --load model_pre.pt --episodes 5000 --opponents heuristic

# Especificar archivo de métricas
python train.py --metrics mi_run.csv
\end{lstlisting}

El entrenamiento guarda automáticamente:
\begin{itemize}[noitemsep]
  \item \texttt{best\_model.pt} — mejor modelo según evaluación formal.
  \item \texttt{model.pt} — checkpoint cada 500 episodios.
  \item \texttt{metrics.csv} — historial de tasas de victoria.
\end{itemize}

\subsection{Evaluación}

\begin{lstlisting}[style=bash]
# Evaluar el mejor modelo (200 partidas vs heurístico y vs aleatorio)
python eval.py

# Evaluar otro checkpoint con más partidas
python eval.py --model model.pt --games 400
\end{lstlisting}

Salida de ejemplo:
\begin{lstlisting}[style=bash]
Resultados vs heuristic (200 partidas):
  Red Neuronal (RED):    98 victorias  (49.0%)
  Heuristico (BLUE):    48 victorias  (24.0%)
  ...
  Objetivo (Fase 3):   46.5%
  Red neuronal:        49.0%  OBJETIVO SUPERADO
\end{lstlisting}

\subsection{Visualización del progreso}

\begin{lstlisting}[style=bash]
# Abrir gráfico interactivo
python plot.py

# Guardar como imagen PNG
python plot.py --save grafico.png

# Usar un archivo de métricas específico
python plot.py --file mi_run.csv
\end{lstlisting}

El gráfico muestra:
\begin{itemize}[noitemsep]
  \item Línea azul: tasa de victoria durante entrenamiento (media móvil, ventana 30).
  \item Puntos rojos: evaluaciones formales contra el heurístico cada 100 episodios.
  \item Línea gris punteada: baseline aleatorio (25\%).
  \item Línea naranja punteada: objetivo de Fase 3 (46.5\%).
\end{itemize}


% ──────────────────────────────────────────────────────────────────────────────
\section{Flujo Completo Recomendado}

\begin{enumerate}
  \item \textbf{Pre-entrenamiento rápido} contra jugadores aleatorios para que la red
        aprenda lo básico (colocar colonias, construir, etc.):
        \begin{lstlisting}[style=bash]
python train.py --episodes 2000 --opponents random --save model_pre.pt
        \end{lstlisting}

  \item \textbf{Entrenamiento principal} contra el heurístico para superar el 46.5\%:
        \begin{lstlisting}[style=bash]
python train.py --load model_pre.pt --episodes 8000 --save model_final.pt
        \end{lstlisting}

  \item \textbf{Monitorear el progreso} con el gráfico (en otra terminal):
        \begin{lstlisting}[style=bash]
python plot.py
        \end{lstlisting}

  \item \textbf{Evaluación final} para comparar con los benchmarks:
        \begin{lstlisting}[style=bash]
python eval.py --model best_model.pt --games 200
        \end{lstlisting}
\end{enumerate}


% ──────────────────────────────────────────────────────────────────────────────
\section{Desafíos Técnicos}

\begin{description}
  \item[Credit assignment] La recompensa llega al final del juego
        (50--100 decisiones). REINFORCE atribuye el mismo peso a todas las
        decisiones, lo cual es una aproximación gruesa. La Fase 4 aborda esto
        con PPO y discounting.

  \item[Información imperfecta] El estado del juego incluye las cartas de todos
        los jugadores (información completa). Una mejora futura es enmascarar
        los recursos de los rivales para simular juego real.

  \item[Espacio de acciones variable] Catan tiene acciones que dependen
        completamente del estado (localización de caminos, colonias, etc.).
        La máscara binaria sobre un espacio fijo de 244 acciones resuelve
        esto de forma eficiente.

  \item[Acción DISCARD] Cuando un jugador tiene más de 7 cartas al caer un 7,
        debe descartar la mitad. Esta acción tiene un espacio combinatorio enorme
        y se maneja con selección aleatoria en la implementación actual.

  \item[Múltiples jugadores] REINFORCE estándar asume un entorno de un solo
        agente. Con 4 jugadores, el entorno es no-estacionario desde la
        perspectiva de la red.
\end{description}


% ──────────────────────────────────────────────────────────────────────────────
\section{Próximos Pasos -- Fase 4}

\begin{itemize}
  \item Implementar \textbf{PPO} (Proximal Policy Optimization) para
        entrenamiento más estable.
  \item \textbf{Self-play} contra versiones anteriores de la red en lugar
        del agente heurístico fijo.
  \item Añadir \textbf{función de valor} (actor-critic) para reducir la
        varianza del gradiente.
  \item Modelar el \textbf{comercio marítimo} en el espacio de acciones de
        forma más completa.
  \item Explorar \textbf{comercio entre jugadores} (el mayor desafío único de Catan).
\end{itemize}

\vfill
\begin{center}
  \rule{0.5\textwidth}{0.4pt}\\[0.3em]
  \small Proyecto Catan AI $\cdot$ Fase 3 completada $\cdot$
  Python + catanatron + PyTorch
\end{center}

\end{document}
